%% ============================================================
%% chapters/conclusion.tex
%% Conclusion — Vedant Patil Master's Thesis
%% ============================================================

\chapter{Conclusion}
\label{ch:conclusion}

%% ============================================================
\section{What This Thesis Set Out To Do}
\label{sec:concl_intro}
%% ============================================================
Far-edge AI systems are increasingly being asked to operate in environments that the tools used to build them were never designed for. Agricultural fields without internet access. Wildlife conservation sites without cellular infrastructure. Physical environments where the standard assumptions of containerized, cloud-native software, that registries are reachable, that network interfaces are stable, that compute resources are abundant, simply do not hold. The question this thesis set out to answer was not whether AI can work at the far edge. It can, and the models are good. The question was why the platforms carrying those models keep failing when they reach the field, and what it would take to build platforms that do not.

%% ============================================================
\section{What Was Found}
\label{sec:concl_rq1}
%% ============================================================
The first deployment, OpenPASS at a soybean farm in Stoutsville, Ohio, answered the first part of that question directly. A system that had passed every lab test failed in three distinct ways within hours of reaching a real agricultural field. The failures shared a common structure: each one traced back to an assumption the platform had made about its operating environment, an assumption that the field did not honor. Enterprise firewalls blocked the pod network. Hardware constraints crashed the inference containers at startup. Zero connectivity made pod reconstruction impossible. None of these were bugs in the application code. They were mismatches between what the infrastructure expected and what the site provided. Naming that category of failure as an inflection point, and building a taxonomy around it, turned what would otherwise be isolated incidents into a classifiable and addressable problem class. The three architectural solutions that resolved each inflection point, VXLAN overlay networking, init container startup sequencing, and air-gapped image pre-population with local Helm repositories, gave OpenPASS the ability to function at Stoutsville without any of the infrastructure it had previously depended on.

The second deployment, the Smart Backpack at Tuttle Park in Columbus, Ohio, answered the second part. Taking the offline-first principles that the Stoutsville experience had forced into existence and applying them proactively, before anything broke, on a Raspberry Pi 4B running Docker Compose instead of an x86 laptop running K3S, in a wildlife monitoring application instead of a digital agriculture one, produced a system that completed 38 autonomous drone missions out of 41 attempts over 20 continuous hours with a 92.7\% success rate and zero network dependency at runtime. The infrastructure failures that Stoutsville encountered were absent at Tuttle Park. Not because Tuttle Park was a friendlier environment, but because the architecture did not give those failures a surface to attach to. No external registry to fail. No physical network interface to firewall. No simultaneous startup spike to crash the containers.

%% ============================================================
\section{The Argument Both Deployments Make Together}
\label{sec:concl_rq2}
%% ============================================================
Taken separately, each deployment proves something limited. Stoutsville proves that far-edge AI systems fail in predictable ways and that those failures can be resolved. Tuttle Park proves that the Smart Backpack worked for 20 hours at a park in Columbus. Taken together, they provide evidence for something more general. The same design principle, treating external connectivity as unavailable from the start rather than as a resource to degrade gracefully from, held across two hardware architectures, two orchestration frameworks, two application domains, and two geographic field sites.

%% P2.3 FIX: Hedge the generalizability claim throughout this section
Our results are consistent with offline-first design principles generalizing across the two orchestration frameworks, hardware architectures, and application domains we tested. K3S and Docker Compose are different systems. x86 and ARM64 are different hardware. Digital agriculture and wildlife monitoring are different problems. In every context where the principle was applied, it produced a system that the field could not break at the infrastructure level. Whether this holds across a broader range of deployment contexts remains to be established through further independent validation.

That is the thesis's central empirical contribution. Offline-first design is not a feature of K3S. It is not a property of containerized agriculture applications. It is a design philosophy for any AI system that must operate where connectivity cannot be assumed, and our evidence is consistent with its validity extending beyond the specific platform or domain it was applied to here.

%% ============================================================
\section{Limitations}
\label{sec:concl_rq3}
%% ============================================================
Three limitations of this work are worth stating precisely.

The Smart Backpack was validated at Tuttle Park using a person as the proxy detection subject rather than animals in a real conservation environment. The pipeline itself is subject-agnostic: what the camera trap sends is a thresholded detection event with GPS coordinates, and the system's response does not vary with what triggered the detection. But the full end-to-end scenario that motivated the design, autonomous aerial follow-up on a wildlife behavioral event at a remote site without available connectivity, has not been tested under those exact conditions. Ungulate behavior in a 220-acre enclosure involves variable movement, vegetation occlusion, GPS multipath, and detection confidence that fluctuates with animal posture and group density, all of which differ meaningfully from a person standing in an open field. The 92.7\% success rate should be interpreted in the context of these proxy conditions. The planned deployment at The Wilds Conservation Center was not completed. That validation remains future work.

The OpenPASS deployment at Stoutsville produced qualitative results. The three inflection points are real, the solutions resolved them, and the system operated successfully after each fix was applied. What is absent is quantitative instrumentation: latency measurements before and after each fix, uptime percentages, resource utilization traces. Grafana and Loki were part of the Chapter 3 observability stack but were not used to capture before-and-after baselines during the Stoutsville deployment. The contributions of this chapter are therefore architectural rather than empirical in the benchmarking sense. They are valid on those terms, but the magnitude of improvement each solution provides cannot be stated with precision.

The inflection-point database proposed in Chapter 3 does not exist yet. Its value as a community artifact is argued from first principles and from the demonstrated cost of each team rediscovering the same deployment failures independently. That argument is sound, but it remains an argument. The database has not been built, populated beyond the three entries this thesis contributes, or evaluated for adoption by practitioners in the far-edge AI community.

%% ============================================================
\section{Future Work}
\label{sec:concl_agnosticism}
%% ============================================================
Three directions follow naturally from where this thesis ends.

The most direct next step is a field deployment of the Smart Backpack at The Wilds Conservation Center or a comparable remote conservation site where no connectivity is available at runtime. Tuttle Park demonstrated that the system works and that its architecture does not depend on connectivity. A deployment at The Wilds would validate the same architecture under the exact environmental conditions that motivated its design, completing the loop between the problem The Wilds posed and the solution Tuttle Park tested. This is not a large research undertaking. The system is built. The site relationship exists. What is needed is the deployment.

The second direction is building the inflection-point database as a community artifact. The three entries this thesis contributes form a starting point: network environment mismatch, resource contention at startup, and runtime registry dependency. The database design proposed in Chapter 3 already specifies the structure each entry should follow: trigger condition, failure mode, architectural resolution, and validation status. What is needed is a hosting mechanism, a contribution protocol, and early adoption by other far-edge AI research groups. The ICICLE NSF AI Institute, where this work is affiliated, provides a natural community for seeding that adoption. A database with ten entries across five deployment sites would already represent a research asset that does not currently exist.

The third direction is multi-drone coordination within the OpenPassLite framework. The drop-in mission architecture of OpenPassLite was designed to make new flight behaviors composable without modifying the framework core. That composability anticipates but does not implement the case where multiple drones operate from a shared ground station, with SmartFields coordinating missions across assets rather than serializing them behind a single threading lock. A wildlife monitoring deployment with two drones and two camera trap networks would double the detection coverage without requiring any changes to the core framework. The threading lock that currently serializes missions would need to become a per-drone scheduling mechanism, and SmartFields would need a simple priority queue. This is a tractable extension of existing work, not a new system.

%% ============================================================
\section{Closing}
\label{sec:concl_limitations}
%% ============================================================
A soybean field in Stoutsville broke a system that the lab said was ready. That failure turned out to be informative. The same class of failure will break every far-edge AI system built on cloud-native assumptions unless those assumptions are examined and eliminated before the system leaves the building. Examining them is not difficult. It requires asking, for every runtime dependency, what happens when this is not available, and then designing so that the answer is: the system keeps running.

Two deployments, different in nearly every dimension, produced results consistent with that answer. That agreement is what this thesis offers: not a prescription for any particular platform or domain, but evidence that the question is worth asking and that asking it early is enough to change the outcome.

\nocite{zhang2020whole,mana2024sustainable,boubin2019managing,suo20233,morris2018model,stewart2004profile,boubin2022marble,morris2016sprint,wang2007autoparam,romero2023ai,barbedo2026explainability,kline2026wildwing,kline2023framework,kline2025studying,kline2024characterizing,dong2025does,bohnett2026understanding,kline2025smartwilds,grushchak2024decentralized,taylor2021avis,xu2016blending,deng2012adaptive,anthropic2024mcp,chen2023frugalgpt,cornell2021birdnet,ong2024routellm,kim2023prism,wildlifeacoustics2023,wani2025smartwilds,kline2024sec,kline2023acsos,kline2024aaai,kline2024acsos,kline2025wildwing,oquinn2024parallelism,kelley2017answer,morris2018sprinting,kelley2016power,boubin2019agriculture,marble2022,hu2021lora,kholiavchenko2024kabr,dwork2006differential,touvron2023llama,abdin2024phi3,shi2024edge,iqbal2024llmplatform,zheng2024judging,mireshghallah2024llmsecret,shao2024privacyassistants,radosevich2025mcp,errico2025mcp,zhang2025agentsecurity,zhan2025prism,nissenbaum2004privacy,ding2024hybridllm,wang2024privatelora}